\section{Methods}


\subsection{Experimental}

Figure \ref{fig:expt_setup} shows the experimental setup. 

HVOF torch with syringe pump potassium feeding. 

Throughout the experiment, the HVOF process is monitored with a variety of diagnostics. 


\begin{figure}[h]
    \includegraphics[width=\linewidth]{\repodir/figures/output/Fig_expt_setup.png} 
    \caption{Experimental setup of the Photoionization experiment. }
    \label{fig:expt_setup}
\end{figure}

\begin{itemize}
\item Determine kr vs nK, position, laser power
\item Absorption spectroscopy for nK
\item Microwave scattering (MWS) for kr 
\end{itemize}


\begin{figure}[p]
    \includegraphics[width=\linewidth]{\repodir/experiment/analysis/various/output/figures/tc_plot_2023-05-18.png} 
    \caption{Representative experiment}
    \label{fig:expt_example}
\end{figure}

A sample experiment is shown in Figure \ref{fig:expt_example}.

There are a variety of independent variables that have been controlled through out the experiments to explore the photoionization process. For the HVOF process, the total mass flow, equivalence ratio, and atomic potassium mass fraction are controlled. Additionally, the position of the mobile MWS and AES diagnostics are controlled, as well as the excimer laser power. 

Experiments consist of a variety of 'test cases', where the the system is brought to a specific combination of the independent variables and measurements are performed. In particular, 'sequences' of test cases are performed where one of these independent variables is changed. 

For the results presented in this report, the equivalence ratio is held constant at a value of 0.8, the total mass flow rate is held constant at 13 g/s. These variables have been varied in supplementary test case sequences, with data shown in the SI. In this report, we focus on three sequences where the potassium mass fraction, the laser power, and the position of the laser downstream from the exit of the free jet are varied. Notably, when not being varied, the position of the AES/MWS diagnostics is held fixed downstream of the torch at a 'nominal position' of 178 mm. The motivation for this position is discussed in the results section. 

Test case sequences (TODO: improve names)
\begin{itemize}
    \item 536 Position: Varying MWS/AES diagnostic position. Full laser power. Kwt = 1\%
    \item 536 Power: Varying laser power. MWS/AES diagnostic at nominal position. Kwt = 1\%
    \item 53x: Varying potassium mass fraction. MWS/AES diagnostic at nominal position. Full laser power.
\end{itemize}



We note that the syringe pump based feeding method has the disadvantage of a finite amount of emulsion that can be fed into the torch. This limits the duration of a given test case sequence, with shorter durations at high potassium mass fraction. For this reason, it was note possible to maintain consistent durations for the different test case sequences. Furthermore, the finite time allowed by the syringe pump means that exploration of experimental conditions was limited. 


Diagnostic processing is shown in Figure \ref{fig:diagnostic_processing}.

\begin{figure}[h]
    \includegraphics[width=\linewidth]{\repodir/figures/output/diagnostic_processing.png} 
    \caption{Processing of Microwave Scattering and Absorption-Emission spectroscopy data}
    \label{fig:diagnostic_processing}
\end{figure}


\subsubsection{Microwave scattering}

 In previous experiments we have determined that the optical emission lifetime is on the order of 10 nanoseconds (SI), which is much shorter than what is predicted for electron-ion recombination. This fast optical decay is consistent with a chemiluminescence process associated with laser-induced dissociation of KOH. [] 

To probe the free electron density, we developed a MWS diagnostic to probe the free electron density and measure the electron-ion recombination rate. Similar diagnostic tools have been developed in the literature.[] The instrument uses a 90 Ghz Gunn microwave source, IQ mixer, etc. The propagation of the radiation through space is not well characterized, so the system measures relative changes instead of absolute transmission. The system measures the change in the transmitted microwave power when a UV pulse passes through the center of the horns. The UV pulse is created with a GAM Excimer laser with 10ns pulse. 

The measurements and processing of the system are shown in Figure \ref{fig:diagnostic_processing}. The IQ mixer is connected to a Lecroy oscilloscope and the in phase (I) and quadrature (Q) signals are recorded with each laser pulse as shown in Figure \ref{fig:diagnostic_processing}A. The magnitude of the transmitted microwaves is shown in \ref{fig:diagnostic_processing}B and is calculated with the following equation.

\begin{equation}
    Mag(t) = \sqrt{I(t)^2 + Q(t)^2}
\end{equation}

The signal is then converted to a transmission by 

\begin{equation}
    T(t) = \frac{Mag(t)}{Mag_{Pre Pulse}}
\end{equation}

Where $Mag_{Pre Pulse}$ is the average magnitude within a time window of one microsecond before the laser pulse. 

Finally, we calculate the combined absorption and scattering, AS, by

\begin{equation}
    AS(t) = 1 - T(t)
\end{equation}

We fit the AS data with a kinetic model. 

\begin{equation}
    \frac{d\Delta n_e}{dt} = - k_r [2 n_{K+,0}\Delta n_e + \Delta{n_e}^2]
\end{equation}

\subsubsection{Absorption spectroscopy}

The absorption emission spectroscopy is used to determine the potassium number density, $n_K$. the AES system uses a 780 nm LED source and a high resolution HR4000 spectrometer. $n_K$ is determined at the exit of the barrel and at the center of the microwave horns (see Figure \ref{fig:expt_setup}).

TODO: how much should this be copied from AIAA 2020

The spectrally dependent absorption, $\alpha(\lambda)$, is measured from the intensity of the LED with no free jet, $I_0$, the intensity
of the seeded jet by itself, $I_F$ , and the intensity of the seeded free jet and LED together $I_{L+F}$


\begin{equation}
    \alpha(\lambda) = 1 - \left(\frac{I_{L+F}(\lambda) - I_{F}(\lambda)}{I_F(\lambda)}\right)
\end{equation}

Currently the absorption is fit with the Beer-Lambert law, but we are working on numerical integration of the absorption to determine the number density and this section will be rewritten. 

The differential equation is 

\begin{equation}
    \frac{dI}{dx} = -\kappa(\lambda) n_K(x) I
\end{equation}

If we assume $n_K$ is constant, we can solve the differential equation to get the Beer lambert law:

\begin{equation}
    \alpha(\lambda) = 1 - \exp(-\kappa(\lambda) n_K L)
\end{equation}

Where $L$ is the path length.

For numerical integration, we write $n_K(x) = n_{K,max} f(x)$, where $f(x)$ is the profile of the number density along the beam path determined by CFD. 




\subsection{Cantera}

We assess the viability of the photoionization process to produce a net return in energy with 0D chemical equilibrium simulations. The methods are described in detail in the SI. Briefly, the modeling is performed with 0D chemical equilibrium calculations to determine the composition and thermophysical properties of a seeded oxy-kerosene combustion products at various temperatures and pressures, approximating various locations in a MHD generator. We then apply a nonequilibrium ionization model to determine the change in electrical conductivity due to photoionization. We then calculate the MHD power output for various photoionization power inputs and determine the regimes in which photoionization is a net energy return. We calculate a figure of merit for the photoionization process, $\alpha_{PI}$, which is the ratio of the change in MHD power output, $dP_{MHD}$ for an intestinal amount of photoionization power input, $dP_{in}$. 


TODO: replace alpha with another symbol to not be confused with absorption coefficient.


\begin{equation}
\alpha_{PI} = \left. \frac{{dP}_{MHD}}{{dP}_{in}} \right|_{P_{in}=0}
\end{equation}


\begin{figure}[h]
    \includegraphics[width=\linewidth]{\repodir/modeling/analysis/output/alpha_curve_demo.png} 
    \caption{Figure of merit for photoionization for various temperatures and pressures. The line}   
    \label{fig:alpha_line_demo}
\end{figure}

Figure \ref{fig:alpha_line_demo} shows $\alpha_{PI}$ for various temperatures and pressures.  For $\alpha_{PI} > 1$ at low temperatures and pressures (blue region), the photoionization process exhibits a net energy return and is viable. For $\alpha_{PI} < 1$, at high temperatures and pressures (red region) the photoionization process is not viable. The line $\alpha_{PI} = 1$ is the boundary between these two regions. To determine the boundary line (green dashed  line), we interpolate $\alpha_{PI}$ each temperature to determine the pressure at which $\alpha_{PI} = 1$. These boundary lines are then shown for various conditions in the results section.

\subsection{CFD}

CFD modeling 

\begin{itemize}
\item nK+ for MWS kr determination
\item CFD results at goldilocks need examination
\item nK vs position profile, for line integral determination of nK
\item Methodology is not developed yet. 
\end{itemize}
